\section{Hard Problems and Where COSI Sits}
\label{sec:hardproblems}

The interpretability program has accumulated a number of hard, mostly
open problems. We list those we take to be load-bearing, characterize
the state of the literature on each, and locate the COSI program inside
them. The list is not exhaustive; it is an attempt to draw the map at
the resolution that makes the COSI contribution legible.

\paragraph{(P1) The faithfulness problem.}
Reverse-engineered circuits and recovered features are explanations of
model behavior, but the gap between an explanation and the actual
mechanism the model uses is real. Wang et al.~\cite{wang2023interpretability}
acknowledge this explicitly for the indirect-object circuit; Conmy et
al.~\cite{conmy2023towards} treat circuit discovery as an automated
search problem precisely because hand-crafted explanations cannot be
trusted to be faithful. \emph{COSI's relation:} a positive cross-model
isomorphism is correspondence evidence. Causal-intervention
verification is Phase~3 of the program (Section~\ref{sec:discussion}).
The faithfulness problem is the reason Phase~3 is required.

\paragraph{(P2) The polysemanticity / superposition problem.}
Single neurons in trained networks typically respond to multiple
unrelated concepts; this is the empirical fact that originally
motivated SAEs. Elhage et al.~\cite{elhage2022superposition} showed
that superposition is the predicted consequence of representing more
features than there are neurons, and that the implicit feature space is
recoverable through sparse dictionary learning. \emph{COSI's relation:}
the operator-theoretic framing reads the superposition phenomenon as a
finite-rank approximation issue: the forward operator's approximately
invariant subspaces have dimension exceeding the network's neuron
count, requiring sparse coding to express. The COSI test does not
depend on the superposition resolution; it operates on the residual
stream's geometry, which contains the superposed features whether or
not we have a clean dictionary.

\paragraph{(P3) The compositional limits problem.}
Lake and Baroni \cite{lake2018generalization} and Hupkes et
al.~\cite{hupkes2020compositionality} documented that neural sequence
models often fail to generalize compositionally in regimes where humans
trivially succeed. Dziri et al.~\cite{dziri2024faith} extended this to
modern transformers, showing systematic failure modes on multi-step
compositional reasoning. \emph{COSI's relation:} our hypothesis is
about the geometry of compositional representations, not about whether
the models reliably perform composition. A positive COSI result is
consistent with both successful and failed compositional generalization;
the relationship between geometric isomorphism and behavioral
compositionality is itself a question we treat as Phase~4 of the
program.

\paragraph{(P4) The cross-model transfer-of-mechanism question.}
If two models implement the same circuits, can interpretation results
on one transfer to the other? The Platonic Representation Hypothesis
\cite{huh2024platonic} and the universality results
\cite{olah2020zoom, chughtai2023toy, gurnee2024universal} suggest yes;
the empirical demonstrations remain limited. Bansal et
al.~\cite{bansal2021revisiting} provide the strongest functional-level
evidence through model stitching; CKA \cite{kornblith2019similarity}
and related measures provide statistical-level evidence. \emph{COSI's
relation:} a positive COSI result, particularly with the Procrustes
transformation $R^*$ explicitly recovered, makes cross-model transfer
of mechanism geometrically tractable. A direction $v_A$ in $V_A$ has a
specific image $R^{*-1} v_A \in V_B$, predicting an intervention site
in $M_B$ that should have the corresponding effect. Phase~3 tests this
prediction.

\paragraph{(P5) The training-distribution confound.}
Two models trained on heavily overlapping web corpora may show
representational similarity for reasons that have nothing to do with
substrate-independent computation: shared training distribution
produces shared statistical residue. Bansal et
al.~\cite{bansal2021revisiting} and Huh et al.~\cite{huh2024platonic}
both flag this as the chief alternative explanation for cross-model
similarity. \emph{COSI's relation:} this is the most important
confound the present calibration does \emph{not} address. Phase~2 of
the program is structured around it (Section~\ref{sec:discussion}).
A positive COSI result that survives a training-distribution swap is
evidence for substrate-independent computation; one that does not is
evidence for shared corpus residue.

\paragraph{(P6) The interpretability scalability problem.}
Recovered circuits and feature dictionaries grow polynomially with
model scale. The interpretability cost per model has so far scaled with
model size, despite the empirical universality findings suggesting the
\emph{computational structure} should scale much more slowly.
\emph{COSI's relation:} if the operator-theoretic framing holds, then
the interpretability of one model in a family transfers via Procrustes
alignment to the other models in the family; the per-model cost is
amortized across the family. This is a research-program-level
implication of COSI, not a calibration-level one.

\paragraph{(P7) The biological--artificial bridge.}
Neural population recordings reveal low-dimensional manifolds
\cite{gallego2017neural, saxena2019towards} whose geometry is preserved
across animals and behaviors. Goal-driven artificial networks predict
these recordings well \cite{yamins2016using}. Whether the same
geometric structure underlies both is an empirical question that has
been gestured at more often than tested directly. \emph{COSI's
relation:} the COSI methodology is, in principle, applicable to
biological--artificial alignment as well, with neural population
activity replacing one model's residual stream. We treat this as
Phase~5; the present paper is restricted to the artificial--artificial
case.

\paragraph{(P8) The substrate-independence question proper.}
The claim that compositional reasoning is substrate-independent in the
strong sense---that the structure of the computation is preserved
under arbitrary changes of substrate, subject to capacity constraints
\cite{ashby1956introduction}---is the deepest claim the COSI program
engages. The free-energy principle \cite{friston2010free} predicts the
claim from variational arguments; cybernetics
\cite{wiener1948cybernetics} and autopoiesis
\cite{maturana1980autopoiesis} provide the structural language for
how substrate-independent computation could be implemented and
maintained. \emph{COSI's relation:} a fully successful program (Phases
0--3 confirming, Phase~4 demonstrating mechanism transfer) would
constitute the sharpest empirical evidence for substrate-independence
that the field has produced for compositional reasoning. We do not
overclaim what Phase~0 alone establishes; the proper claim is the
calibration result reported here, with the path forward specified.

These eight problems are not independent. The faithfulness problem (P1)
is the precondition for the cross-model transfer question (P4). The
training-distribution confound (P5) is the principal alternative
explanation for any positive COSI result. The substrate-independence
question (P8) is the eventual stake. The COSI program is structured
around them: Phase~0 calibrates the infrastructure (this paper),
Phase~1 amplifies architectural distance, Phase~2 addresses (P5),
Phase~3 addresses (P1) and (P4) jointly, and the integration with (P7)
and (P8) is the subject of the synthesis paper that this program is
building toward.
